% %--------------------------
% %	CONFIGURATION
% %--------------------------
\documentclass[11pt,a4paper]{moderncv}
\moderncvstyle{classic}
\moderncvcolor{black}
\definecolor{firstnamecolor}{rgb}{0,0,0}

% Basic packages
\usepackage[utf8]{inputenc}
\usepackage[T1]{fontenc}
\usepackage[scale=0.85]{geometry}
\usepackage[dvipsnames]{xcolor}
\definecolor{DarkBlue}{RGB}{0, 0, 150}
\usepackage[colorlinks=true, urlcolor=DarkBlue, linkcolor=DarkBlue]{hyperref}

\usepackage{microtype}

% Configure font settings
\renewcommand{\rmdefault}{ppl}
\renewcommand{\sfdefault}{phv}
\renewcommand*{\namefont}{\fontsize{26}{18}\mdseries}

% Configure microtype
\microtypesetup{expansion=false}

% Load sensitive information
\input{../../other_material/personal_info_dk.tex}

% Title document
\title{Curriculum Vitae}

% %--------------------------

\begin{document}
\makecvtitle

\section*{Introduction}
I am a pragmatic data scientist and ML engineer with seven years of experience building, deploying, and maintaining machine learning systems across pharmaceutical analytics, healthcare, manufacturing, and e-commerce. I have owned production ML pipelines, modernised forecasting systems, deployed ML products in Databricks, and worked across data science, software engineering, and stakeholder-facing product delivery. I am especially interested in Corti's Senior MLOps Engineer role because it combines reliability and clinically meaningful AI systems.

% ---------------------------
% Selected Work Experience
% ---------------------------
\section{Selected Work Experience}

% ---------------------------
\large{Full-Stack Data Scientist \& Project Manager @ Signum Life Science}
\normalsize
\begin{itemize}
    \item Architecting and building an end-to-end forecasting platform for pharmaceutical price and demand. Includes ETL pipelines, model training workflows, MLflow experiment tracking and versioning to support thousands of simultaneous in-production models. All built with Databricks asset bundles and IaC.
    \item Owned an AWS-hosted next-best-action tool. Developed XAI features together with UX researchers.
\end{itemize}

{\footnotesize\textit{Keywords: Databricks, IaC, MLflow, Docker, AWS, XGBoost, forecasting, explainable AI, pharmaceutical analytics}}

\hfill{\small{\textit{\hyperref[sec:signum]{Read more...}}}}

% ---------------------------
\large{Principal Data Scientist @ Valcon}
\normalsize
\begin{itemize}
    \item Architected and built an end-to-end, real-time, production system for natural language inputs.
    \item Mentored colleagues and led trainings in explainable AI and Git.
    \item Awarded "Valcon People's Choice Award" for exemplary work, mentoring, team spirit, and curiosity.
\end{itemize}

{\footnotesize\textit{Keywords: production ML, neural networks, NLP, ML engineering, mentoring}}

\hfill{\small{\textit{\hyperref[sec:valcon]{Read more...}}}}

% ---------------------------
\large{Researcher within Data Science @ RaySearch Laboratories}
\normalsize
\begin{itemize}
    \item Built ML methods for automated radiotherapy planning from segmented 3D CT scans.
    \item Used autoencoders, sparse Gaussian processes, and optimisation for dose planning.
    \item Collaborated with clinicians and medical physicists on clinical workflow fit.
\end{itemize}

{\footnotesize\textit{Keywords: medical AI, clinical decision support, 3D CT, autoencoders, probabilistic modelling}}

\hfill{\small{\textit{\hyperref[sec:raysearch]{Read more...}}}}

\section{Education}

\cventry{2015--2020}{Royal Institute of Technology, KTH -- Data Science and Engineering}{}{GPA: 5.0 (out of 5.0)}{}{Master: Statistical Learning and Data Analytics, Bachelor: Engineering Physics. \href{https://kth.diva-portal.org/smash/get/diva2:1431327/FULLTEXT02.pdf}{See thesis}. \newline Relevant coursework: Statistical Machine Learning; Optimization; Applied Time Series.}

\cventry{2019--2019}{Swiss Federal Institute of Technology, ETH -- Applied Mathematics}{Exchange}{}{}{}

\cventry{2017--2022}{Stockholm University, SU -- Economics}{Supplementary Bachelor's Degree}{}{}{}

% ---------------------------
% Skills and Other
% ---------------------------
\section{Skills and Other}
\normalsize
\cvitem{Technology}{Python, SQL, PyTorch, TensorFlow, NLP, Computer Vision, MLFlow, IaC, Databricks, AWS, Azure, FastAPI, MongoDB, Qdrant, ETL pipelines.}
\cvitem{Soft-skills}{Project management, mentoring, training, stakeholder communication, commercialisation support.}
\cvitem{Language}{\textit{Native}: English, Swedish. \textit{Intermediate}: Danish.}

\newpage

% ---------------------------
% Detailed Work Experience
% ---------------------------
\section{Detailed Work Experience}

% ---------------------------
\phantomsection\label{sec:signum}
\cventry{2024--Present}{Full-Stack Data Scientist and Project Manager}{Signum Life Science}{Copenhagen}{}{
    At Signum, my current work is on an end-to-end Databricks forecasting platform for pharmaceutical price and demand forecasting. I work across architecture, project management, data science, ML engineering, and software engineering for the platform. The work builds on a pharmaceutical price forecasting model that helps clients navigate Denmark's fortnightly blind auctions. After stabilising the XGBoost-based system, I am now driving efforts toward a more robust and modern implementation.
    \newline\hspace*{2em}
    Earlier in the role, I owned and extended Kwarts, a next-best-action engine for pharmaceutical sales representatives. Originally built externally, the system was extended with features such as explainable AI in collaboration with UX designers, and supported commercialisation by aiding our sales representatives. Kwarts adapts to varied customer data maturity, from rule-based logic to deep learning, and is deployed in AWS for scalability.
    \newline\hspace*{2em}
    Beyond model development, I organised an internal week-long hackathon that included up to 40 participants to foster innovation and collaboration between teams.
}{}

% ---------------------------
\phantomsection\label{sec:valcon}
\cventry{2022--2024}{Principal Data Scientist}{Valcon}{Copenhagen}{}{
    Led data science projects across pharma, energy, IoT, and manufacturing. Mentored junior colleagues in software architecture and data science best practices, led explainable AI and Git trainings, and regularly presented complex technical topics to non-technical stakeholders.
    \newline\hspace*{2em}
    One project involved training a real-time deep neural network with a custom training pipeline and modular architecture. The training data consisted of natural-language inputs and high-dimensional multi-class labels, which made for a challenging prediction problem. I was the lead data scientist and architect of the solution and responsible for implementing the end-to-end production pipeline.
    \newline\hspace*{2em}
    I was awarded "Valcon People's Choice Award" for exemplary work, mentoring, team spirit, and curiosity.
}{}

% ---------------------------
\phantomsection\label{sec:valtech}
\cventry{2021--2022}{Data Scientist}{Valtech}{Copenhagen}{}{
    Delivered ML solutions in e-commerce. Built time series forecasts using Facebook Prophet, and applied clustering to behavioural data to support personalisation. Worked closely with business stakeholders to ensure actionable outputs.
}{}

% ---------------------------
\phantomsection\label{sec:raysearch}
\cventry{2019--2021}{Researcher within Data Science}{RaySearch Laboratories}{Stockholm}{}{
    At RaySearch Laboratories, I worked on automating radiotherapy planning using a combination of deep learning, probabilistic modelling, and multi-objective optimisation. I developed a pipeline that extracted latent anatomical features from segmented 3D CT scans using autoencoders, and used these representations to identify similar patients via a learned similarity metric. Dose-volume histograms and spatial dose distributions were then predicted using sparse Gaussian processes, providing a probabilistic framework for dose planning conditioned on patient geometry.
    \newline\hspace*{2em}
    Alongside this, I redesigned some of RaySearch's lexicographic optimisation algorithms to better reflect clinical decision hierarchies, particularly in palliative care where trade-offs between tumour control and organ sparing are nuanced. I collaborated closely with medical physicists and clinicians throughout, ensuring the models aligned with both oncological standards and practical workflow requirements.
}{}

% ---------------------------
% Other Work Experience
% ---------------------------
\section{Other Detailed Experience}

% ---------------------------
\phantomsection\label{sec:iris}
\cventry{2025--2025}{Solo Project}{Project Iris}{}{}{
    In this personal project I built a PoC end-to-end computer vision system for making product images on e-commerce sites interactive. The pipeline detects individual items using a YOLOv8 model, and embeds them semantically using OpenCLIP. The embeddings are then compared to support intelligent product linking across all images and products in a web store.
    \newline\hspace*{2em}
    The system includes a full-stack setup with a FastAPI backend, MongoDB for persistence, a Qdrant index for vector search, and a custom-built Chromium plug-in that overlays clickable links directly onto web shop imagery in real time. The project also includes a dynamic scraping and indexing pipeline for ingesting and structuring product data from arbitrary e-commerce sites.
    \newline\hspace*{2em}
    The project was a learning opportunity for fine-tuning of deep learning models and practical ML engineering.
}{}

\end{document}
